% ============================================================================
\chapter{模群 $\SL_2(\Z)$}
\label{ch:modular-group}
% ============================================================================

% ----------------------------------------------------------------------------
\section{上半平面与模群的作用}
% ----------------------------------------------------------------------------

\begin{definition}[上半平面]
\textbf{上半平面}定义为
\[
  \HH = \{ \tau \in \C \mid \im(\tau) > 0 \}.
\]
\end{definition}

\begin{definition}[模群]
\textbf{模群} $\SL_2(\Z)$ 是所有行列式为 $1$ 的 $2 \times 2$ 整数矩阵组成的群：
\[
  \SL_2(\Z) = \left\{
    \begin{pmatrix} a & b \\ c & d \end{pmatrix}
    \;\middle|\; a, b, c, d \in \Z,\; ad - bc = 1
  \right\}.
\]
它通过\textbf{莫比乌斯变换}作用在 $\HH$ 上：
\[
  \gamma \cdot \tau = \frac{a\tau + b}{c\tau + d},
  \quad \gamma = \begin{pmatrix} a & b \\ c & d \end{pmatrix}.
\]
\end{definition}

% ----------------------------------------------------------------------------
\section{生成元 $S$ 和 $T$}
% ----------------------------------------------------------------------------

% TODO: S = ((0,-1),(1,0)), T = ((1,1),(0,1))，证明它们生成 SL₂(ℤ)

% ----------------------------------------------------------------------------
\section{基本域}
% ----------------------------------------------------------------------------

% TODO: 标准基本域 F 的描述与证明

% ----------------------------------------------------------------------------
\section{尖点}
% ----------------------------------------------------------------------------

% TODO: 扩展上半平面 H* = H ∪ Q ∪ {∞}，尖点的定义

% ----------------------------------------------------------------------------
\section{习题}
% ----------------------------------------------------------------------------
